\chapter{Conceptual Model}\label{chap:conceptual_model}

The literature reviewed in the previous chapter identifies a set of factors that together define the design space for the proposed system. This chapter narrows that space to the specific model adopted in this thesis, providing the reasoning that connects each design decision to the evidence base.

\section{Key Factors and Their Relative Importance}

Five factors emerge from the literature as determinants of the system design, though they do not carry equal weight.

The first two factors establish why machine learning is necessary at all. \textbf{Multi-variable mould risk} means that no single threshold on any single sensor can reliably represent spoilage hazard --- temperature, humidity, and VOC concentration interact, and only a multi-input classifier can model that relationship. \textbf{Early VOC signal} means that the sensor types required must include gas sensing alongside temperature and humidity: the volatile compounds produced during early \textit{Botrytis cinerea} activity provide a leading indicator of mould onset before any visible decay is established \cite{Yang2025, Ma2023}. For \textit{B. cinerea} specifically, detectable biomarkers in the storage atmosphere of infected strawberry include 1-octen-3-ol, 2-methyl-1-butanol, 2-methyl-1-propanol, and short-chain ketones \cite{Vandendriessche2012}. These two factors together fix the sensor selection --- DHT22 for temperature and humidity, SGP30 for total VOC, MQ3 for ethanol vapour --- and justify the choice of a neural network classifier over a rule-based approach. Beyond the instantaneous sensor readings, the rate of change of VOC concentration ($\Delta$TVOC) was also included as a feature: a rising VOC level is more indicative of active and accelerating fungal metabolism than a static reading alone, and including the delta allows the classifier to distinguish between a stable elevated-VOC environment and one in which concentrations are actively increasing.

The third factor, \textbf{sensor placement below produce level}, is a physical constraint that follows from the stratification behaviour of the target VOC compounds in low-airflow enclosed environments \cite{Zhang2026}. It does not alter the classifier design but is a necessary condition for valid signal capture, and it informs the mechanical design of the sensor node described in the following chapter.

The final two factors carry the most weight in shaping the research questions. \textbf{Connectivity independence} eliminates cloud-based inference as an option for vehicles in transit and rural storage facilities where network access cannot be assumed, satisfying functional requirement FR6. \textbf{Concept drift} eliminates a centrally pre-trained static model as a reliable long-term solution, since the statistical relationship between sensor readings and mould risk shifts between deployment environments and changes over time \cite{Bayram2022, Ilic2023}. Together, these two factors force the question the thesis addresses: if the model must run and adapt locally on a battery-powered node, what does that cost in energy, and at what point on the adaptation spectrum is that cost operationally justified?

\section{Mapping the Model to Research Questions and Requirements}

The proposed system is a battery-powered sensor node --- an ESP32 microcontroller with DHT22, SGP30, and MQ3 sensors, running a neural network classifier entirely on-device without cloud connectivity. This single architecture is evaluated under three configurations corresponding to the three TinyML frameworks: TF Lite Micro (inference only, no on-device adaptation), TinyOL (partial adaptation of the final output layer), and AIfES (full on-device training with no pre-trained weights). The progression from inference-only to full training maps directly onto the progression from connectivity dependence to full autonomy identified in the literature. Figure~\ref{fig:conceptual_model} illustrates this architecture.

\begin{figure}[htbp]
\centering
\begin{tikzpicture}[
    sensor/.style={draw, rectangle, rounded corners=3pt, minimum width=2.6cm, minimum height=0.8cm, fill=gray!12, font=\small, align=center},
    nnbox/.style={draw, rectangle, rounded corners=3pt, minimum width=3.4cm, minimum height=2.2cm, fill=blue!12, font=\scriptsize, align=center},
    outbox/.style={draw, rectangle, rounded corners=3pt, minimum width=2.6cm, minimum height=1.2cm, fill=green!12, font=\small, align=center},
    arrow/.style={-Stealth, semithick}
]
% Sensor nodes (left column)
\node[sensor] (dht) at (0,  1.6) {DHT22\\\scriptsize Temp / Humidity};
\node[sensor] (sgp) at (0,  0)   {SGP30\\\scriptsize Total VOC};
\node[sensor] (mq3) at (0, -1.6) {MQ3\\\scriptsize Ethanol};
% ESP32 bounding box
\draw[draw=black!50, fill=blue!4, rounded corners=5pt] (2.7,-2.3) rectangle (6.3, 2.3);
\node[font=\footnotesize\bfseries] at (4.5, 2.0) {ESP32};
% NN classifier inside ESP32
\node[nnbox] (nn) at (4.5, -0.1) {Neural Network\\Classifier\\[5pt]\scriptsize TF Lite Micro (frozen)\\TinyOL (last layer)\\AIfES (full training)};
% Output node (right)
\node[outbox] (out) at (8.6, 0) {Mould Risk\\[3pt]\scriptsize Low~/~High};
% Sensor to ESP32 arrows
\draw[arrow] (dht.east) -- ++(0.5,0) |- (2.7,  1.1);
\draw[arrow] (sgp.east) -- (2.7, 0);
\draw[arrow] (mq3.east) -- ++(0.5,0) |- (2.7, -1.1);
% ESP32 to output arrow
\draw[arrow] (6.3, 0) -- (out.west);
\end{tikzpicture}
\caption{Architecture of the autonomous mould detection node. Sensor readings are passed to the ESP32, which runs one of three TinyML classifier configurations depending on the experimental condition. No external connectivity is required at any stage.}
\label{fig:conceptual_model}
\end{figure}

This architecture is designed to answer the primary research question --- what is the energy cost of full on-device training compared to partial adaptation and inference-only execution, and when is that cost operationally justified --- by providing three directly comparable implementations on identical hardware. The four sub-questions each address a distinct layer of that answer. The first isolates the energy cost of the sensing hardware alone, so that the overhead added by machine learning can be attributed specifically to the framework rather than to the underlying sensors. The second characterises the energy consumed per training cycle and per inference cycle across all three framework variants, producing the per-operation figures needed for comparison. The third verifies that the three variants achieve comparable prediction performance, since an energy comparison between frameworks of substantially different accuracy would not be meaningful. The fourth combines the measured energy figures with a realistic deployment duty cycle to project how long a battery-powered node could operate under each framework --- translating the per-operation measurements into an operational judgement.

The cloud-free, fully autonomous operating requirement is not an incidental constraint --- it is the condition that makes the energy question meaningful in the first place. An architecture with reliable network access could offload training to a remote server entirely, and the on-device energy cost would be irrelevant. It is precisely the removal of that option, motivated by the connectivity and concept-drift factors identified above, that creates the need to understand what on-device training costs and whether a battery-powered node can sustain it over a realistic operational period.

\section{Summary}

The conceptual model is an autonomous, battery-powered sensor node that classifies mould risk from continuous environmental measurements without any dependency on cloud connectivity or manual recalibration. The choice of sensor types follows from the identified VOC biomarkers; the sensor placement follows from the stratification behaviour of those compounds; the multi-input classifier follows from the multi-variable nature of mould risk; and the spectrum of TinyML frameworks follows from the requirement to operate and adapt locally under concept drift. The energy cost of each point on that spectrum is the question the experimental work addresses.
